\documentclass{article}

% if you need to pass options to natbib, use, e.g.:
%     \PassOptionsToPackage{numbers, compress}{natbib}
% before loading aaj2026


% ready for submission
\usepackage[final]{aaj2026}
\bibliographystyle{plainnat}
\usepackage[numbers]{natbib}

% to compile a preprint version, e.g., for submission to arXiv, add add the
% [preprint] option:
%     \usepackage[preprint]{aaj2026}


% to compile a camera-ready version, add the [final] option, e.g.:
%     \usepackage[final]{aaj2026}


% to avoid loading the natbib package, add option nonatbib:
%    \usepackage[nonatbib]{aaj2026}


\usepackage[utf8]{inputenc} % allow utf-8 input
\usepackage[T1]{fontenc}    % use 8-bit T1 fonts
\usepackage[hyphens]{url}
\usepackage{hyperref}       % hyperlinks
\usepackage{url}            % simple URL typesetting
\usepackage{booktabs}       % professional-quality tables
\usepackage{amsfonts}       % blackboard math symbols
\usepackage{nicefrac}       % compact symbols for 1/2, etc.
\usepackage{microtype}      % microtypography
\usepackage{xcolor}         % colors


\title{World Models for the Electricity Grid \\ (v 2026.07.04)}


% The \author macro works with any number of authors. There are two commands
% used to separate the names and addresses of multiple authors: \And and \AND.
%
% Using \And between authors leaves it to LaTeX to determine where to break the
% lines. Using \AND forces a line break at that point. So, if LaTeX puts 3 of 4
% authors names on the first line, and the last on the second line, try using
% \AND instead of \And before the third author name.


\author{%
  Nathan Murthy\thanks{As of this writing, Nathan serves as the Head of Engineering at Verse; however, in no way does this position paper reflect the opinions of the company or discuss any of its proprietary technology} \\
  Verse \\
  San Francisco, CA \\
  \texttt{nathan.murthy@\{gmail.com,verse.inc\}} \\
  % examples of more authors
% \And
  % Jane S.~Hippocampus\\
  % Department of Computer Science\\
  % AI Agent Journal\\
  % Singapore \\
  % \texttt{jane@aiagentjournal.org} \\
  % \AND
  % Coauthor \\
  % Affiliation \\
  % Address \\
  % \texttt{email} \\
  % \And
  % Coauthor \\
  % Affiliation \\
  % Address \\
  % \texttt{email} \\
  % \And
  % Coauthor \\
  % Affiliation \\
  % Address \\
  % \texttt{email} \\
}

% remove "1st AI Agent Journal (2026)" text in the footer of
% the first page of `\usepackage[final]{aaj2026}`
\makeatletter
\def\@noticestring{}
\def\aft@noticestring{}
\makeatother

\begin{document}


\maketitle


\begin{abstract}
  The bulk electric power grid---the largest machine ever built on the planet---is undergoing a
  once-in-a-generation transformation as demand for accelerated compute to power artificial
  intelligence (AI) workloads has put inexorable capacity strain on available generation,
  transmission, and distribution systems. Solar, wind, and battery technologies will be deployed
  en masse to alleviate AI interconnection bottlenecks; in turn, optimal dispatch of these
  resources to support large AI workloads will also benefit from new AI models. World model
  architectures have the potential to enhance coarse-grain model-predictive control (MPC) and
  reinforcement learning (RL) techniques used in industry by providing differentiable, 
  action-conditioned predictors over latent representations of grid state---offering a tractable
  approximation to the grid's dynamics without expensive grid simulators.  In particular,
  Joint-Embedding Predictive Architecture (JEPA) can be applied to hierarchical planning and 
  operation of these grid systems toward highly autonomous grid intelligence and thus maximally
  aware, efficient, reliable, and sustainable grid management.

\end{abstract}


\section{Background: industry strain today}

On the eve commemorating the 250th anniversary of the United States' founding, the U.S.
Department of Energy (DoE) issued its forty-fifth emergency order of the year to electric grid 
operators ~\cite{Frazin2026} ~\cite{DoE_202c_2026}---an all-time record barely halfway through
2026 more than double the previous year---as temperature readings also broke records in over
280 locations with nearly half the country placed under extreme heat warnings by the National 
Weather Service ~\cite{Ittai2026}. Just 9 months earlier, Federal Energy Regulatory Commission
(FERC) Chairman David Rosner expressed an urgency shared by many in industry and government
that “America needs every  electron and every molecule of every type we can get, and we need
more infrastructure to move it" ~\cite{Howland2025}. This is a belief shared not only in the
United States but in other major geopolitical regions: Chinese Premier Li Qiang emphasized
efforts to advance grid construction and accelerate energy infrastructure optimization
~\cite{Xinua2026} and European think tanks have recommended policymakers to look at alleviating
interconnection bottlenecks with lead times much longer than in the US~\cite{Hajdari2026}. The
execution requirements and the administrative labor which follows for businesses and regulators
alone may even dwarf the complexity of just physically building new systems ~\cite{Sguazzi2026}.

Fundamentally, each power grid stakeholder needs tools
for capacity planning, connection studies, and optimal dispatch; many of these tools today suffer
from a patchwork of various solutions that all model the grid and its dynamics from short to long
time scales without the kind of integrated intelligence needed to mitigate the gamut of planning
and operational concerns. The solution to these problems is both an engineering and coordination 
challenge ~\cite{Grindal2026}, and AI will be a key element of realizing integrated solutions as
much as it has been a catalyst for revealing where these stressors exist today.

Several companies have deployed solutions to meet these challenges over the years \cite{OldEnergySoftware}, and many new 
companies have sprouted in recent years that have approached these problems with fresher eyes 
~\cite{Sguazzi2026} ~\cite{Bessemer2026}.
Each of these solutions continue to rely on one critical decision-making input: \emph{forecasts}.
Companies either develop their own internal coarse-grained forecasts or vendor them from third
parties with very limited visibility into the internals of how the forecasts are generated and/or 
with prohibitive complexity and cost of proprietary grid modeling and simulation software.
Moreover, once a forecast has been generated and decision made based on that forecast, there is
very limited and even unknown feedback pertaining to the accuracy of the forecast and the 
causal dynamics of the decision made. Human decision makers are overwhelmed with information
and data coming from the grid on multiple time frames and are thus severely handicapped in their
ability to make optimal decisions under uncertainty compared to agentic systems with full system
context and capabilities to reason about the consequences of their actions in both latent and
real time frames. Because of this, businesses and governments would benefit tremendously from 
world models for and of the electricity grid. The scope, complexity, and impact of such models,
I would argue, eclipse even the scientific progress in AI research applied to fields such as
protein folding and biochemistry for sheer ubiquity of electrical energy in the modern age for
virtually all aspects of life.

The dominant techniques employed for these decision-making forecasts span classical statistical
methods (e.g. ARIMA, Monte Carlo simulation), MPC (in situations where dispatch is required), and
more recently RL in the academic literature. Capacity expansion and market operation modeling
often involve solving a mixed-integer linear program (MILP) with tools like PLEXOS and
optimization libraries such as Gurobi, CBC, or cvxpy, at each step of a recurring dynamic program.
One challenge in industry today is that because many of these models are proprietary, it is
difficult to truly know the quality, accuracy, and validity of the techniques and their outputs
without spending large sums on licenses and access to bulky models---the competitive drive among 
businesses who profit from grid planning and operations have very little incentive or desire to 
disseminate their methods in any truly scientific manner, and so sophisticated, high-quality
forecasts may be indistinguishable from simple, low-quality forecasts except for price point. This
presents a significant opportunity to go beyond the industry status quo, especially as grid
planning, operations, and studies become increasingly strained with each new data center
interconnection request. The status quo has not moved much beyond traditional CPU-based
computational methods. Newer forecasting tools and AI reasoning capabilities built on accelerated
computing technologies will offer several advantages over current forecasting methods.


\subsection{Style}


Papers to be submitted to AI Agent Journal 2026 must be prepared according to the
instructions presented here. Papers may only be up to {\bf nine} pages long,
including figures. Additional pages \emph{containing only acknowledgments and
references} are allowed. Papers that exceed the page limit will not be
reviewed, or in any other way considered for presentation at the conference.


The margins in 2026 are the same as those in previous years.


Authors are required to use the AI Agent Journal 2026 \LaTeX{} style files obtainable at the
AI Agent Journal 2026 website as indicated below. Please make sure you use the current files
and not previous versions. Tweaking the style files may be grounds for
rejection.


\subsection{Retrieval of style files}


The style files for AI Agent Journal 2026 and other conference information are available on
the website at
\begin{center}
  \url{https://www.aiagentjournal.org/}
\end{center}
The file \verb+aaj2026.pdf+ contains these instructions and illustrates the
various formatting requirements your AI Agent Journal 2026 paper must satisfy.


The only supported style file for AI Agent Journal 2026 is \verb+aaj2026.sty+,
rewritten for \LaTeXe{}.  \textbf{Previous style files for \LaTeX{} 2.09,
  Microsoft Word, and RTF are no longer supported!}


The \LaTeX{} style file contains three optional arguments: \verb+final+, which
creates a camera-ready copy, \verb+preprint+, which creates a preprint for
submission to, e.g., arXiv, and \verb+nonatbib+, which will not load the
\verb+natbib+ package for you in case of package clash.


\paragraph{Preprint option}
If you wish to post a preprint of your work online, e.g., on arXiv, using the
AI Agent Journal 2026 style, please use the \verb+preprint+ option. This will create a
nonanonymized version of your work with the text ``Preprint. Work in progress.''
in the footer. This version may be distributed as you see fit, as long as you do not say which conference it was submitted to. Please \textbf{do
  not} use the \verb+final+ option, which should \textbf{only} be used for
papers accepted to AI Agent Journal 2026.


At submission time, please omit the \verb+final+ and \verb+preprint+
options. This will anonymize your submission and add line numbers to aid
review. Please do \emph{not} refer to these line numbers in your paper as they
will be removed during generation of camera-ready copies.


The file \verb+aaj2026.tex+ may be used as a ``shell'' for writing your
paper. All you have to do is replace the author, title, abstract, and text of
the paper with your own.


The formatting instructions contained in these style files are summarized in
Sections \ref{gen_inst}, \ref{headings}, and \ref{others} below.


\section{General formatting instructions}
\label{gen_inst}


The text must be confined within a rectangle 5.5~inches (33~picas) wide and
9~inches (54~picas) long. The left margin is 1.5~inch (9~picas).  Use 10~point
type with a vertical spacing (leading) of 11~points.  Times New Roman is the
preferred typeface throughout, and will be selected for you by default.
Paragraphs are separated by \nicefrac{1}{2}~line space (5.5 points), with no
indentation.


The paper title should be 17~point, initial caps/lower case, bold, centered
between two horizontal rules. The top rule should be 4~points thick and the
bottom rule should be 1~point thick. Allow \nicefrac{1}{4}~inch space above and
below the title to rules. All pages should start at 1~inch (6~picas) from the
top of the page.


For the final version, authors' names are set in boldface, and each name is
centered above the corresponding address. The lead author's name is to be listed
first (left-most), and the co-authors' names (if different address) are set to
follow. If there is only one co-author, list both author and co-author side by
side.


Please pay special attention to the instructions in Section \ref{others}
regarding figures, tables, acknowledgments, and references.


\section{Headings: first level}
\label{headings}


All headings should be lower case (except for first word and proper nouns),
flush left, and bold.


First-level headings should be in 12-point type.


\subsection{Headings: second level}


Second-level headings should be in 10-point type.


\subsubsection{Headings: third level}


Third-level headings should be in 10-point type.


\paragraph{Paragraphs}


There is also a \verb+\paragraph+ command available, which sets the heading in
bold, flush left, and inline with the text, with the heading followed by 1\,em
of space.


\section{Citations, figures, tables, references}
\label{others}


These instructions apply to everyone.


\subsection{Citations within the text}


The \verb+natbib+ package will be loaded for you by default.  Citations may be
numeric or author/year, as long as you maintain internal consistency. As to the
format of the references themselves, any style is acceptable as long as it is
used consistently. Suggested way is to use~\cite{Alexander1995}, or this for multiple articles~\cite{Bower1995,Hasselmo1995}.


The documentation for \verb+natbib+ may be found at
\begin{center}
  \url{http://mirrors.ctan.org/macros/latex/contrib/natbib/natnotes.pdf}
\end{center}
Of note is the command \verb+\citet+, which produces citations appropriate for
use in inline text.  For example,
\begin{verbatim}
   \citet{hasselmo} investigated\dots
\end{verbatim}
produces
\begin{quote}
  Hasselmo, et al.\ (1995) investigated\dots
\end{quote}


If you wish to load the \verb+natbib+ package with options, you may add the
following before loading the \verb+aaj2026+ package:
\begin{verbatim}
   \PassOptionsToPackage{options}{natbib}
\end{verbatim}


If \verb+natbib+ clashes with another package you load, you can add the optional
argument \verb+nonatbib+ when loading the style file:
\begin{verbatim}
   \usepackage[nonatbib]{aaj2026}
\end{verbatim}



\subsection{Footnotes}


Footnotes should be used sparingly.  If you do require a footnote, indicate
footnotes with a number\footnote{Sample of the first footnote.} in the
text. Place the footnotes at the bottom of the page on which they appear.
Precede the footnote with a horizontal rule of 2~inches (12~picas).


Note that footnotes are properly typeset \emph{after} punctuation
marks.\footnote{As in this example.}


\subsection{Figures}


\begin{figure}
  \centering
  \fbox{\rule[-.5cm]{0cm}{4cm} \rule[-.5cm]{4cm}{0cm}}
  \caption{Sample figure caption.}
\end{figure}


All artwork must be neat, clean, and legible. Lines should be dark enough for
purposes of reproduction. The figure number and caption always appear after the
figure. Place one line space before the figure caption and one line space after
the figure. The figure caption should be lower case (except for first word and
proper nouns); figures are numbered consecutively.


You may use color figures.  However, it is best for the figure captions and the
paper body to be legible if the paper is printed in either black/white or in
color.


\subsection{Tables}


All tables must be centered, neat, clean and legible.  The table number and
title always appear before the table.  See Table~\ref{sample-table}.


Place one line space before the table title, one line space after the
table title, and one line space after the table. The table title must
be lower case (except for first word and proper nouns); tables are
numbered consecutively.


Note that publication-quality tables \emph{do not contain vertical rules.} We
strongly suggest the use of the \verb+booktabs+ package, which allows for
typesetting high-quality, professional tables:
\begin{center}
  \url{https://www.ctan.org/pkg/booktabs}
\end{center}
This package was used to typeset Table~\ref{sample-table}.


\begin{table}
  \caption{Sample table title}
  \label{sample-table}
  \centering
  \begin{tabular}{lll}
    \toprule
    \multicolumn{2}{c}{Part}                   \\
    \cmidrule(r){1-2}
    Name     & Description     & Size ($\mu$m) \\
    \midrule
    Dendrite & Input terminal  & $\sim$100     \\
    Axon     & Output terminal & $\sim$10      \\
    Soma     & Cell body       & up to $10^6$  \\
    \bottomrule
  \end{tabular}
\end{table}

\subsection{Math}
Note that display math in bare TeX commands will not create correct line numbers for submission. Please use LaTeX (or AMSTeX) commands for unnumbered display math. (You really shouldn't be using \$\$ anyway; see \url{https://tex.stackexchange.com/questions/503/why-is-preferable-to} and \url{https://tex.stackexchange.com/questions/40492/what-are-the-differences-between-align-equation-and-displaymath} for more information.)

\subsection{Final instructions}

Do not change any aspects of the formatting parameters in the style files.  In
particular, do not modify the width or length of the rectangle the text should
fit into, and do not change font sizes (except perhaps in the
\textbf{References} section; see below). Please note that pages should be
numbered.


\section{Preparing PDF files}


Please prepare submission files with paper size ``US Letter,'' and not, for
example, ``A4.''


Fonts were the main cause of problems in the past years. Your PDF file must only
contain Type 1 or Embedded TrueType fonts. Here are a few instructions to
achieve this.


\begin{itemize}


\item You should directly generate PDF files using \verb+pdflatex+.


\item You can check which fonts a PDF files uses.  In Acrobat Reader, select the
  menu Files$>$Document Properties$>$Fonts and select Show All Fonts. You can
  also use the program \verb+pdffonts+ which comes with \verb+xpdf+ and is
  available out-of-the-box on most Linux machines.


\item \verb+xfig+ "patterned" shapes are implemented with bitmap fonts.  Use
  "solid" shapes instead.


\item The \verb+\bbold+ package almost always uses bitmap fonts.  You should use
  the equivalent AMS Fonts:
\begin{verbatim}
   \usepackage{amsfonts}
\end{verbatim}
followed by, e.g., \verb+\mathbb{R}+, \verb+\mathbb{N}+, or \verb+\mathbb{C}+
for $\mathbb{R}$, $\mathbb{N}$ or $\mathbb{C}$.  You can also use the following
workaround for reals, natural and complex:
\begin{verbatim}
   \newcommand{\RR}{I\!\!R} %real numbers
   \newcommand{\Nat}{I\!\!N} %natural numbers
   \newcommand{\CC}{I\!\!\!\!C} %complex numbers
\end{verbatim}
Note that \verb+amsfonts+ is automatically loaded by the \verb+amssymb+ package.


\end{itemize}


If your file contains type 3 fonts or non embedded TrueType fonts, we will ask
you to fix it.


\subsection{Margins in \LaTeX{}}


Most of the margin problems come from figures positioned by hand using
\verb+\special+ or other commands. We suggest using the command
\verb+\includegraphics+ from the \verb+graphicx+ package. Always specify the
figure width as a multiple of the line width as in the example below:
\begin{verbatim}
   \usepackage[pdftex]{graphicx} ...
   \includegraphics[width=0.8\linewidth]{myfile.pdf}
\end{verbatim}
See Section 4.4 in the graphics bundle documentation
(\url{http://mirrors.ctan.org/macros/latex/required/graphics/grfguide.pdf})


A number of width problems arise when \LaTeX{} cannot properly hyphenate a
line. Please give LaTeX hyphenation hints using the \verb+\-+ command when
necessary.

\begin{ack}
Use unnumbered first level headings for the acknowledgments. All acknowledgments
go at the end of the paper before the list of references. Moreover, you are required to declare
funding (financial activities supporting the submitted work) and competing interests (related financial activities outside the submitted work).
More information about this disclosure can be found at: \url{https://aiagentjournal.org}.


Do {\bf not} include this section in the anonymized submission, only in the final paper. You can use the \texttt{ack} environment provided in the style file to automatically hide this section in the anonymized submission.
\end{ack}

\section*{Notes for References}


References follow the acknowledgments in the camera-ready paper. Use unnumbered first-level heading for
the references. Any choice of citation style is acceptable as long as you are
consistent. It is permissible to reduce the font size to \verb+small+ (9 point)
when listing the references.
Note that the Reference section does not count towards the page limit.
\medskip


\bibliography{biblio.bib}


%%%%%%%%%%%%%%%%%%%%%%%%%%%%%%%%%%%%%%%%%%%%%%%%%%%%%%%%%%%%

\appendix

\section{Appendix / supplemental material}


Optionally include supplemental material (complete proofs, additional experiments and plots) in appendix.
All such materials \textbf{SHOULD be included in the main submission.}

%%%%%%%%%%%%%%%%%%%%%%%%%%%%%%%%%%%%%%%%%%%%%%%%%%%%%%%%%%%%




\end{document}
